

\begin{table}[H]
    \begin{tabular}{|llllllll|}
    \hline
    \multicolumn{1}{|l|}{Caso de Uso}       & \multicolumn{5}{l|}{Registrar pago}                                                                          & \multicolumn{2}{l|}{\cellcolor[HTML]{E6E6E6}CU-19}          \\ \hline
    \multicolumn{1}{|l|}{Actores}           & \multicolumn{7}{l|}{{\color[HTML]{808080} \textit{Empleado, Administrador}}}                                                                                                                           \\ \hline
    \multicolumn{1}{|l|}{Tipo}              & \multicolumn{7}{l|}{{\color[HTML]{808080} \textit{Básico, Esencial}}}                                                                                                                \\ \hline
    \multicolumn{1}{|l|}{Referencias}       & \multicolumn{2}{l|}{{\color[HTML]{808080} \textit{RF-3}}}              & \multicolumn{5}{l|}{}                                                                                       \\ \hline
    \multicolumn{1}{|l|}{Precondición}      & \multicolumn{7}{l|}{{\color[HTML]{808080} \textit{Existen una reserva y un empleado que pueda registrar el pago.}}}                                                                                                \\ \hline
    \multicolumn{1}{|l|}{Poscondición}      & \multicolumn{7}{l|}{}                                                                                                                                                                  \\ \hline
    \multicolumn{1}{|l|}{Autor}             & \multicolumn{1}{l|}{}                        & \multicolumn{2}{l|}{Fecha}           & \multicolumn{1}{l|}{}     & \multicolumn{2}{l|}{Versión}         & 1.0                           \\ \hline
    \multicolumn{8}{|l|}{\cellcolor[HTML]{F2F2F2}Propósito}                                                                                                                                                                          \\ \hline
    \multicolumn{8}{|p{15cm}|}{{\color[HTML]{808080} \textit{Registrar el pago de una reserva.}}}                                                                                                           \\ \hline
    \multicolumn{8}{|l|}{\cellcolor[HTML]{F2F2F2}Resumen}                                                                                                                                                                            \\ \hline
    \multicolumn{8}{|p{15cm}|}{{\color[HTML]{808080} \textit{El empleado se encargará de, una vez se haya producido una reserva, registrar el pago correspondiente a la misma.}}} \\ \hline
    \end{tabular}
\end{table}

\begin{table}[H]
    \begin{tabular}{|llllllll|}
    \hline
    \multicolumn{1}{|l|}{Caso de Uso}       & \multicolumn{5}{l|}{Emitir factura}                                                                          & \multicolumn{2}{l|}{\cellcolor[HTML]{E6E6E6}CU-20}          \\ \hline
    \multicolumn{1}{|l|}{Actores}           & \multicolumn{7}{l|}{{\color[HTML]{808080} \textit{Empleado, Administrador}}}                                                                                                                           \\ \hline
    \multicolumn{1}{|l|}{Tipo}              & \multicolumn{7}{l|}{{\color[HTML]{808080} \textit{Básico, Esencial}}}                                                                                                                \\ \hline
    \multicolumn{1}{|l|}{Referencias}       & \multicolumn{2}{l|}{{\color[HTML]{808080} \textit{RF-3.2}}}              & \multicolumn{5}{l|}{}                                                                                       \\ \hline
    \multicolumn{1}{|l|}{Precondición}      & \multicolumn{7}{l|}{{\color[HTML]{808080} \textit{Existe una reserva a la que asociar una factura.}}}                                                                                                \\ \hline
    \multicolumn{1}{|l|}{Poscondición}      & \multicolumn{7}{l|}{}                                                                                                                                                                  \\ \hline
    \multicolumn{1}{|l|}{Autor}             & \multicolumn{1}{l|}{}                        & \multicolumn{2}{l|}{Fecha}           & \multicolumn{1}{l|}{}     & \multicolumn{2}{l|}{Versión}         & 1.0                           \\ \hline
    \multicolumn{8}{|l|}{\cellcolor[HTML]{F2F2F2}Propósito}                                                                                                                                                                          \\ \hline
    \multicolumn{8}{|p{15cm}|}{{\color[HTML]{808080} \textit{Registrar el pago de una reserva.}}}                                                                                                           \\ \hline
    \multicolumn{8}{|l|}{\cellcolor[HTML]{F2F2F2}Resumen}                                                                                                                                                                            \\ \hline
    \multicolumn{8}{|p{15cm}|}{{\color[HTML]{808080} \textit{Una vez se haga una reserva, se debe generar una factura para dejar constancia del pago de cara a la contabilización de los flujos monetarios de la empresa.}}} \\ \hline
    \end{tabular}
\end{table}

\begin{table}[H]
    \begin{tabular}{|llllllll|}
    \hline
    \multicolumn{1}{|l|}{Caso de Uso}       & \multicolumn{5}{l|}{Enviar factura a cliente}                                                                          & \multicolumn{2}{l|}{\cellcolor[HTML]{E6E6E6}CU-21}          \\ \hline
    \multicolumn{1}{|l|}{Actores}           & \multicolumn{7}{l|}{Empleado, Administrador}                                                                                                                           \\ \hline
    \multicolumn{1}{|l|}{Tipo}              & \multicolumn{7}{l|}{Primario, Esencial}                                                                                                                \\ \hline
    \multicolumn{1}{|l|}{Referencias}       & \multicolumn{2}{l|}{RF-3.2}              & \multicolumn{5}{l|}{CU-20}                                                                                       \\ \hline
    \multicolumn{1}{|l|}{Precondición}      & \multicolumn{7}{l|}{Existe una factura asociada al cliente que la solicita.}                                                                                                \\ \hline
    \multicolumn{1}{|l|}{Poscondición}      & \multicolumn{7}{l|}{}                                                                                                                                                                  \\ \hline
    \multicolumn{1}{|l|}{Autor}             & \multicolumn{1}{l|}{}                        & \multicolumn{2}{l|}{Fecha}           & \multicolumn{1}{l|}{}     & \multicolumn{2}{l|}{Versión}         & 1.0                           \\ \hline
    \multicolumn{8}{|l|}{\cellcolor[HTML]{F2F2F2}Propósito}                                                                                                                                                                          \\ \hline
    \multicolumn{8}{|p{15cm}|}{Enviar a un cliente la factura asociada a su reserva.}                                                                                                           \\ \hline
    \multicolumn{8}{|l|}{\cellcolor[HTML]{F2F2F2}Resumen}                                                                                                                                                                            \\ \hline
    \multicolumn{8}{|p{15cm}|}{Una vez se haga una reserva, existe la posibilidad de enviar la factura de la reserva al cliente para que tenga constancia de la misma.} \\ \hline
    \end{tabular}
\end{table}

\begin{table}[H]
    \begin{tabular}{|llllllll|}
    \hline
    \multicolumn{1}{|l|}{Caso de Uso}       & \multicolumn{5}{l|}{Descargar Factura}                                                                          & \multicolumn{2}{l|}{\cellcolor[HTML]{E6E6E6}CU-22}          \\ \hline
    \multicolumn{1}{|l|}{Actores}           & \multicolumn{7}{l|}{Empleado, Administrador, Cliente}                                                                                                                           \\ \hline
    \multicolumn{1}{|l|}{Tipo}              & \multicolumn{7}{l|}{Primario, Esencial}                                                                                                                \\ \hline
    \multicolumn{1}{|l|}{Referencias}       & \multicolumn{2}{l|}{RF-3.2}              & \multicolumn{5}{l|}{CU-20}                                                                                       \\ \hline
    \multicolumn{1}{|l|}{Precondición}      & \multicolumn{7}{l|}{Existe una factura que poder descargar.}                                                                                                \\ \hline
    \multicolumn{1}{|l|}{Poscondición}      & \multicolumn{7}{l|}{}                                                                                                                                                                  \\ \hline
    \multicolumn{1}{|l|}{Autor}             & \multicolumn{1}{l|}{}                        & \multicolumn{2}{l|}{Fecha}           & \multicolumn{1}{l|}{}     & \multicolumn{2}{l|}{Versión}         & 1.0                           \\ \hline
    \multicolumn{8}{|l|}{\cellcolor[HTML]{F2F2F2}Propósito}                                                                                                                                                                          \\ \hline
    \multicolumn{8}{|p{15cm}|}{Descargar la factura de una reserva.}                                                                                                           \\ \hline
    \multicolumn{8}{|l|}{\cellcolor[HTML]{F2F2F2}Resumen}                                                                                                                                                                            \\ \hline
    \multicolumn{8}{|p{15cm}|}{Una vez se realice una reserva, tanto el empleado, como el administrador como el cliente tendrán la posibilidad de descargar la factura asociada a la misma.} \\ \hline
    \end{tabular}
\end{table}

\begin{table}[H]
    \begin{tabular}{|llllllll|}
    \hline
    \multicolumn{1}{|l|}{Caso de Uso}       & \multicolumn{5}{l|}{Elegir método de pago}                                                                          & \multicolumn{2}{l|}{\cellcolor[HTML]{E6E6E6}CU-23}          \\ \hline
    \multicolumn{1}{|l|}{Actores}           & \multicolumn{7}{l|}{Cliente}                                                                                                                           \\ \hline
    \multicolumn{1}{|l|}{Tipo}              & \multicolumn{7}{l|}{Primario, Esencial}                                                                                                                \\ \hline
    \multicolumn{1}{|l|}{Referencias}       & \multicolumn{2}{l|}{RF-3.1}              & \multicolumn{5}{l|}{}                                                                                       \\ \hline
    \multicolumn{1}{|l|}{Precondición}      & \multicolumn{7}{l|}{}                                                                                                \\ \hline
    \multicolumn{1}{|l|}{Poscondición}      & \multicolumn{7}{l|}{}                                                                                                                                                                  \\ \hline
    \multicolumn{1}{|l|}{Autor}             & \multicolumn{1}{l|}{}                        & \multicolumn{2}{l|}{Fecha}           & \multicolumn{1}{l|}{}     & \multicolumn{2}{l|}{Versión}         & 1.0                           \\ \hline
    \multicolumn{8}{|l|}{\cellcolor[HTML]{F2F2F2}Propósito}                                                                                                                                                                          \\ \hline
    \multicolumn{8}{|p{15cm}|}{Poder variar la forma de pago de una reserva.}                                                                                                           \\ \hline
    \multicolumn{8}{|l|}{\cellcolor[HTML]{F2F2F2}Resumen}                                                                                                                                                                            \\ \hline
    \multicolumn{8}{|p{15cm}|}{El cliente tendrá la posibilidad de elegir distintos métodos de pago en la web para abonar la cuantía de la reserva.} \\ \hline
    \end{tabular}
\end{table}

\begin{table}[H]
    \begin{tabular}{|llllllll|}
    \hline
    \multicolumn{1}{|l|}{Caso de Uso}       & \multicolumn{5}{l|}{Procesar Reembolsos}                                                                          & \multicolumn{2}{l|}{\cellcolor[HTML]{E6E6E6}CU-24}          \\ \hline
    \multicolumn{1}{|l|}{Actores}           & \multicolumn{7}{l|}{Empleado, Administrador}                                                                                                                           \\ \hline
    \multicolumn{1}{|l|}{Tipo}              & \multicolumn{7}{l|}{Primario, Esencial}                                                                                                                \\ \hline
    \multicolumn{1}{|l|}{Referencias}       & \multicolumn{2}{l|}{RF-3.3}              & \multicolumn{5}{l|}{CU-25, CU-26}                                                                                       \\ \hline
    \multicolumn{1}{|l|}{Precondición}      & \multicolumn{7}{l|}{Existe una reserva previa de la que procesar un reembolso.}                                                                                                \\ \hline
    \multicolumn{1}{|l|}{Poscondición}      & \multicolumn{7}{l|}{}                                                                                                                                                                  \\ \hline
    \multicolumn{1}{|l|}{Autor}             & \multicolumn{1}{l|}{}                        & \multicolumn{2}{l|}{Fecha}           & \multicolumn{1}{l|}{}     & \multicolumn{2}{l|}{Versión}         & 1.0                           \\ \hline
    \multicolumn{8}{|l|}{\cellcolor[HTML]{F2F2F2}Propósito}                                                                                                                                                                          \\ \hline
    \multicolumn{8}{|p{15cm}|}{Llevar a cabo el reembolso de una reserva.}                                                                                                           \\ \hline
    \multicolumn{8}{|l|}{\cellcolor[HTML]{F2F2F2}Resumen}                                                                                                                                                                            \\ \hline
    \multicolumn{8}{|p{15cm}|}{Si un cliente no está de acuerdo con algún aspecto de su reserva, podrá solicitar un reembolso, el cual será procesado bien por un empleado o por el administrador.} \\ \hline
    \end{tabular}
\end{table}

\begin{table}[H]
    \begin{tabular}{|llllllll|}
    \hline
    \multicolumn{1}{|l|}{Caso de Uso}       & \multicolumn{5}{l|}{Consultar Datos Cliente}                                                                          & \multicolumn{2}{l|}{\cellcolor[HTML]{E6E6E6}CU-25}          \\ \hline
    \multicolumn{1}{|l|}{Actores}           & \multicolumn{7}{l|}{Empleado, Administrador, Cliente}                                                                                                                           \\ \hline
    \multicolumn{1}{|l|}{Tipo}              & \multicolumn{7}{l|}{Primario, Esencial}                                                                                                                \\ \hline
    \multicolumn{1}{|l|}{Referencias}       & \multicolumn{2}{l|}{RF-3}              & \multicolumn{5}{l|}{}                                                                                       \\ \hline
    \multicolumn{1}{|l|}{Precondición}      & \multicolumn{7}{l|}{}                                                                                                \\ \hline
    \multicolumn{1}{|l|}{Poscondición}      & \multicolumn{7}{l|}{}                                                                                                                                                                  \\ \hline
    \multicolumn{1}{|l|}{Autor}             & \multicolumn{1}{l|}{}                        & \multicolumn{2}{l|}{Fecha}           & \multicolumn{1}{l|}{}     & \multicolumn{2}{l|}{Versión}         & 1.0                           \\ \hline
    \multicolumn{8}{|l|}{\cellcolor[HTML]{F2F2F2}Propósito}                                                                                                                                                                          \\ \hline
    \multicolumn{8}{|p{15cm}|}{Acceder a los datos de un determinado cliente.}                                                                                                           \\ \hline
    \multicolumn{8}{|l|}{\cellcolor[HTML]{F2F2F2}Resumen}                                                                                                                                                                            \\ \hline
    \multicolumn{8}{|p{15cm}|}{Para realizar diferentes tareas, es necesario acceder a los datos de un cliente para poder obtener otra información del mismo que haga falta.} \\ \hline
    \end{tabular}
\end{table}

\begin{table}[H]
    \begin{tabular}{|llllllll|}
    \hline
    \multicolumn{1}{|l|}{Caso de Uso}       & \multicolumn{5}{l|}{Modificar Reserva}                                                                          & \multicolumn{2}{l|}{\cellcolor[HTML]{E6E6E6}CU-26}          \\ \hline
    \multicolumn{1}{|l|}{Actores}           & \multicolumn{7}{l|}{Empleado, Administrador}                                                                                                                           \\ \hline
    \multicolumn{1}{|l|}{Tipo}              & \multicolumn{7}{l|}{Primario, Esencial}                                                                                                                \\ \hline
    \multicolumn{1}{|l|}{Referencias}       & \multicolumn{2}{l|}{RF-3.3}              & \multicolumn{5}{l|}{}                                                                                       \\ \hline
    \multicolumn{1}{|l|}{Precondición}      & \multicolumn{7}{l|}{Existencia de un reembolso que requiera modificar una reserva.}                                                                                                \\ \hline
    \multicolumn{1}{|l|}{Poscondición}      & \multicolumn{7}{l|}{}                                                                                                                                                                  \\ \hline
    \multicolumn{1}{|l|}{Autor}             & \multicolumn{1}{l|}{}                        & \multicolumn{2}{l|}{Fecha}           & \multicolumn{1}{l|}{}     & \multicolumn{2}{l|}{Versión}         & 1.0                           \\ \hline
    \multicolumn{8}{|l|}{\cellcolor[HTML]{F2F2F2}Propósito}                                                                                                                                                                          \\ \hline
    \multicolumn{8}{|p{15cm}|}{Modificar los datos asociados a una reserva.}                                                                                                           \\ \hline
    \multicolumn{8}{|l|}{\cellcolor[HTML]{F2F2F2}Resumen}                                                                                                                                                                            \\ \hline
    \multicolumn{8}{|p{15cm}|}{En caso de realizarse un reembolso, será necesario que un empleado o administrador modifique la reserva previamente realizada para poder o bien eliminarla o bien modificar algunos aspectos de la misma.} \\ \hline
    \end{tabular}
\end{table}

\begin{table}[H]
    \begin{tabular}{|llllllll|}
    \hline
    \multicolumn{1}{|l|}{Caso de Uso}       & \multicolumn{5}{l|}{Consultar Historial de Pagos}                                                                          & \multicolumn{2}{l|}{\cellcolor[HTML]{E6E6E6}CU-27}          \\ \hline
    \multicolumn{1}{|l|}{Actores}           & \multicolumn{7}{l|}{Administrador, Cliente}                                                                                                                           \\ \hline
    \multicolumn{1}{|l|}{Tipo}              & \multicolumn{7}{l|}{Primario, Esencial}                                                                                                                \\ \hline
    \multicolumn{1}{|l|}{Referencias}       & \multicolumn{2}{l|}{RF-3}              & \multicolumn{5}{l|}{CU-25}                                                                                       \\ \hline
    \multicolumn{1}{|l|}{Precondición}      & \multicolumn{7}{l|}{}                                                                                                \\ \hline
    \multicolumn{1}{|l|}{Poscondición}      & \multicolumn{7}{l|}{}                                                                                                                                                                  \\ \hline
    \multicolumn{1}{|l|}{Autor}             & \multicolumn{1}{l|}{}                        & \multicolumn{2}{l|}{Fecha}           & \multicolumn{1}{l|}{}     & \multicolumn{2}{l|}{Versión}         & 1.0                           \\ \hline
    \multicolumn{8}{|l|}{\cellcolor[HTML]{F2F2F2}Propósito}                                                                                                                                                                          \\ \hline
    \multicolumn{8}{|p{15cm}|}{Ver los pagos realizados por un cliente.}                                                                                                           \\ \hline
    \multicolumn{8}{|l|}{\cellcolor[HTML]{F2F2F2}Resumen}                                                                                                                                                                            \\ \hline
    \multicolumn{8}{|p{15cm}|}{Tanto el administrador (por ser información confidencial) como el cliente podrán acceder al historial de pagos del cliente en cuestión.} \\ \hline
    \end{tabular}
\end{table}

\begin{table}[H]
    \begin{tabular}{|llllllll|}
    \hline
    \multicolumn{1}{|l|}{Caso de Uso}       & \multicolumn{5}{l|}{Configurar Métodos de Pago}                                                                          & \multicolumn{2}{l|}{\cellcolor[HTML]{E6E6E6}CU-28}          \\ \hline
    \multicolumn{1}{|l|}{Actores}           & \multicolumn{7}{l|}{Administrador}                                                                                                                           \\ \hline
    \multicolumn{1}{|l|}{Tipo}              & \multicolumn{7}{l|}{Primario, Esencial}                                                                                                                \\ \hline
    \multicolumn{1}{|l|}{Referencias}       & \multicolumn{2}{l|}{RF-3.1}              & \multicolumn{5}{l|}{}                                                                                       \\ \hline
    \multicolumn{1}{|l|}{Precondición}      & \multicolumn{7}{l|}{}                                                                                                \\ \hline
    \multicolumn{1}{|l|}{Poscondición}      & \multicolumn{7}{l|}{}                                                                                                                                                                  \\ \hline
    \multicolumn{1}{|l|}{Autor}             & \multicolumn{1}{l|}{}                        & \multicolumn{2}{l|}{Fecha}           & \multicolumn{1}{l|}{}     & \multicolumn{2}{l|}{Versión}         & 1.0                           \\ \hline
    \multicolumn{8}{|l|}{\cellcolor[HTML]{F2F2F2}Propósito}                                                                                                                                                                          \\ \hline
    \multicolumn{8}{|p{15cm}|}{Configurar el correcto funcionamiento de los métodos de pago.}                                                                                                           \\ \hline
    \multicolumn{8}{|l|}{\cellcolor[HTML]{F2F2F2}Resumen}                                                                                                                                                                            \\ \hline
    \multicolumn{8}{|p{15cm}|}{El administrador deberá configurar y supervisar a nivel interno que todos los métodos de pago disponibles en la web funcionen de forma correcta además de solucionar los posibles errores producidos en los mismos.} \\ \hline
    \end{tabular}
\end{table}


% Please add the following required packages to your document preamble:
% \usepackage[table,xcdraw]{xcolor}
% Beamer presentation requires \usepackage{colortbl} instead of \usepackage[table,xcdraw]{xcolor}
\begin{table}[H]
    \begin{tabular}{|llllllll|}
    \hline
    \multicolumn{1}{|l|}{\cellcolor[HTML]{F2F2F2}Caso de Uso}           & \multicolumn{5}{l|}{Implementar chat en vivo y asistencia}                                                                       & \multicolumn{2}{l|}{\cellcolor[HTML]{E6E6E6}CU-29}             \\ \hline
    \multicolumn{1}{|l|}{{\color[HTML]{808080} \textit{Actores}}}       & \multicolumn{7}{l|}{{\color[HTML]{808080} \textit{Administrador, Guía Turístico, Empleado}}}                                                                                                      \\ \hline
    \multicolumn{1}{|l|}{Tipo}                                          & \multicolumn{7}{l|}{{\color[HTML]{808080} \textit{Primario, Esencial}}}                                                                                                                           \\ \hline
    \multicolumn{1}{|l|}{Referencias}                                   & \multicolumn{2}{l|}{{\color[HTML]{808080} \textit{RF-4.2}}}                & \multicolumn{5}{l|}{}                                                                                                \\ \hline
    \multicolumn{1}{|l|}{Precondición}                                  & \multicolumn{7}{l|}{{\color[HTML]{808080} \textit{Suponemos que está implementado con IA.}}}                                                                                                      \\ \hline
    \multicolumn{1}{|l|}{Poscondición}                                  & \multicolumn{7}{l|}{}                                                                                                                                                                             \\ \hline
    \multicolumn{1}{|l|}{Autor}                                         & \multicolumn{1}{l|}{}                         & \multicolumn{2}{l|}{Fecha}             & \multicolumn{1}{l|}{}        & \multicolumn{2}{l|}{Versión}             & 1.0                            \\ \hline
    \multicolumn{8}{|l|}{Propósito}                                                                                                                                                                                                                                         \\ \hline
    \multicolumn{8}{|p{15cm}|}{{\color[HTML]{808080} \textit{Usar un chat en vivo para dar soporte y asistencia}}}                                                                                                                                                                                       \\ \hline
    \multicolumn{8}{|l|}{Resumen}                                                                                                                                                                                                                                           \\ \hline
    \multicolumn{8}{|p{15cm}|}{{\color[HTML]{808080} \textit{Este chat está implementado con IA, por lo que actúa de manera autónoma. Cabe destacar que puede fallar, por ende, debemos de tener un empleado que esté para determinadas situaciones que requieran un trato más específico y acertado.}}} \\ \hline
    \end{tabular}
    \end{table}


% Please add the following required packages to your document preamble:
% \usepackage[table,xcdraw]{xcolor}
% Beamer presentation requires \usepackage{colortbl} instead of \usepackage[table,xcdraw]{xcolor}
\begin{table}[H]
    \begin{tabular}{|llllllll|}
    \hline
    \multicolumn{1}{|l|}{Caso de Uso}  & \multicolumn{5}{l|}{Ofrecer sección de preguntas frecuentes}                                       & \multicolumn{2}{l|}{\cellcolor[HTML]{E6E6E6}CU-30}  \\ \hline
    \multicolumn{1}{|l|}{Actores}      & \multicolumn{7}{l|}{{\color[HTML]{808080} \textit{Administrador, Cliente}}}                                                                              \\ \hline
    \multicolumn{1}{|l|}{Tipo}         & \multicolumn{7}{l|}{{\color[HTML]{808080} \textit{Primario, Esencial}}}                                                                                  \\ \hline
    \multicolumn{1}{|l|}{Referencias}  & \multicolumn{2}{l|}{{\color[HTML]{808080} \textit{RF-4.3}}}    & \multicolumn{5}{l|}{}                                                                   \\ \hline
    \multicolumn{1}{|l|}{Precondición} & \multicolumn{7}{l|}{{\color[HTML]{808080} \textit{Los usuarios tengan acceso al foro de preguntas frecuentes.}}}                                         \\ \hline
    \multicolumn{1}{|l|}{Poscondición} & \multicolumn{7}{l|}{}                                                                                                                                    \\ \hline
    \multicolumn{1}{|l|}{Autor}        & \multicolumn{1}{l|}{}                   & \multicolumn{2}{l|}{Fecha} & \multicolumn{1}{l|}{}  & \multicolumn{2}{l|}{Versión} & 1.0                       \\ \hline
    \multicolumn{8}{|l|}{\cellcolor[HTML]{F2F2F2}Propósito}                                                                                                                              \\ \hline
    \multicolumn{8}{|p{15cm}|}{{\color[HTML]{808080} \textit{Dar acceso a un foro de preguntas frecuentes para que los usuarios puedan ver las preguntas más frecuentes.}}}                             \\ \hline
    \multicolumn{8}{|l|}{\cellcolor[HTML]{F2F2F2}Resumen}                                                                                                                                         \\ \hline
    \multicolumn{8}{|p{15cm}|}{{\color[HTML]{808080} \textit{Este foro está destinado a resolver las dudas más frecuentes de los distintos usuarios.}}}                                                 \\ \hline
    \end{tabular}
    \end{table}


% Please add the following required packages to your document preamble:
% \usepackage[table,xcdraw]{xcolor}
% Beamer presentation requires \usepackage{colortbl} instead of \usepackage[table,xcdraw]{xcolor}
\begin{table}[H]
    \begin{tabular}{|llllllll|}
    \hline
    \multicolumn{1}{|l|}{Caso de Uso}       & \multicolumn{5}{l|}{Reservar y Comprar paquetes}                                                                         & \multicolumn{2}{l|}{\cellcolor[HTML]{E6E6E6}CU-31}          \\ \hline
    \multicolumn{1}{|l|}{Actores}           & \multicolumn{7}{l|}{{\color[HTML]{808080} \textit{Cliente}}}                                                                                                                           \\ \hline
    \multicolumn{1}{|l|}{Tipo}              & \multicolumn{7}{l|}{{\color[HTML]{808080} \textit{Primario, Esencial}}}                                                                                                                \\ \hline
    \multicolumn{1}{|l|}{Referencias}       & \multicolumn{2}{l|}{{\color[HTML]{808080} \textit{RF-4.1}}}              & \multicolumn{5}{l|}{{\color[HTML]{808080} \textit{CU-36, CU-32, CU-33}}}                                    \\ \hline
    \multicolumn{1}{|l|}{Precondición}      & \multicolumn{7}{l|}{{\color[HTML]{808080} \textit{Tenga la opción disponible en web.}}}                                                                                                \\ \hline
    \multicolumn{1}{|l|}{Poscondición}      & \multicolumn{7}{l|}{}                                                                                                                                                                  \\ \hline
    \multicolumn{1}{|l|}{Autor}             & \multicolumn{1}{l|}{}                        & \multicolumn{2}{l|}{Fecha}           & \multicolumn{1}{l|}{}     & \multicolumn{2}{l|}{Versión}         & 1.0                           \\ \hline
    \multicolumn{8}{|l|}{\cellcolor[HTML]{F2F2F2}Propósito}                                                                                                                                                                          \\ \hline
    \multicolumn{8}{|p{15cm}|}{{\color[HTML]{808080} \textit{Reservar y comprar paquetes de manera online a través de la web.}}}                                                                                                           \\ \hline
    \multicolumn{8}{|l|}{\cellcolor[HTML]{F2F2F2}Resumen}                                                                                                                                                                            \\ \hline
    \multicolumn{8}{|p{15cm}|}{{\color[HTML]{808080} \textit{Se permite reservar y comprar paquetes de manera online por parte del cliente gracias a la implementación de la web por parte de los trabajadores de la empresa turística.}}} \\ \hline
    \end{tabular}
    \end{table}


% Please add the following required packages to your document preamble:
% \usepackage[table,xcdraw]{xcolor}
% Beamer presentation requires \usepackage{colortbl} instead of \usepackage[table,xcdraw]{xcolor}
\begin{table}[H]
    \begin{tabular}{|llllllll|}
    \hline
    \multicolumn{1}{|l|}{Caso de Uso}                     & \multicolumn{5}{l|}{Información de destinos}                                                                                                                                                & \multicolumn{2}{l|}{\cellcolor[HTML]{E6E6E6}CU-32}                                      \\ \hline
    \multicolumn{1}{|l|}{Actores}                         & \multicolumn{7}{l|}{{\color[HTML]{808080} \textit{Cliente}}}                                                                                                                                                                                                                          \\ \hline
    \multicolumn{1}{|l|}{Tipo}                            & \multicolumn{7}{l|}{{\color[HTML]{808080} \textit{Primario, Esencial}}}                                                                                                                                                                                                               \\ \hline
    \multicolumn{1}{|l|}{Referencias}                     & \multicolumn{2}{l|}{{\color[HTML]{808080} \textit{RF-4.3}}}                                        & \multicolumn{5}{l|}{}                                                                                                                                                            \\ \hline
    \multicolumn{1}{|l|}{Precondición}                    & \multicolumn{7}{l|}{{\color[HTML]{808080} \textit{Estén los destinos registrados con su información.}}}                                                                                                                                                                               \\ \hline
    \multicolumn{1}{|l|}{Poscondición}                    & \multicolumn{7}{l|}{}                                                                                                                                                                                                                                                                 \\ \hline
    \multicolumn{1}{|l|}{Autor}                           & \multicolumn{1}{l|}{}                                     & \multicolumn{2}{l|}{Fecha}                                     & \multicolumn{1}{l|}{}                   & \multicolumn{2}{l|}{Versión}                                     & 1.0                                         \\ \hline
    \multicolumn{8}{|l|}{\cellcolor[HTML]{F2F2F2}Propósito}                                                                                                                                                                                                                                                                                       \\ \hline
    \multicolumn{8}{|p{15cm}|}{{\color[HTML]{808080} \textit{Permitir a los usuarios acceder a un historial de reservas realizadas.}}}                                                                                                                                                                                                                  \\ \hline
    \multicolumn{8}{|l|}{\cellcolor[HTML]{F2F2F2}Resumen}                                                                                                                                                                                                                                                                                         \\ \hline
    \multicolumn{8}{|p{15cm}|}{{\color[HTML]{999999} \textit{El sistema almacena y muestra el historial de reservas realizadas, permitiendo a los usuarios consultar detalles, realizar modificaciones o cancelaciones según las políticas establecidas. Puede incluir filtros y herramientas de búsqueda para facilitar el acceso a la información.}}} \\ \hline
    \end{tabular}
    \end{table}

% Please add the following required packages to your document preamble:
% \usepackage[table,xcdraw]{xcolor}
% Beamer presentation requires \usepackage{colortbl} instead of \usepackage[table,xcdraw]{xcolor}
% Please add the following required packages to your document preamble:
% \usepackage[table,xcdraw]{xcolor}
% Beamer presentation requires \usepackage{colortbl} instead of \usepackage[table,xcdraw]{xcolor}
\begin{table}[H]
    \begin{tabular}{|llllllll|}
    \hline
    \multicolumn{1}{|l|}{Caso de Uso}        & \multicolumn{5}{l|}{Historial de Reservas}                                                                                     & \multicolumn{2}{l|}{\cellcolor[HTML]{E6E6E6}CU-33}            \\ \hline
    \multicolumn{1}{|l|}{Actores}            & \multicolumn{7}{l|}{{\color[HTML]{808080} \textit{Administrador, Guía Turístico, Empleado}}}                                                                                                   \\ \hline
    \multicolumn{1}{|l|}{Tipo}               & \multicolumn{7}{l|}{{\color[HTML]{808080} \textit{Primario, Esencial}}}                                                                                                                        \\ \hline
    \multicolumn{1}{|l|}{Referencias}        & \multicolumn{2}{l|}{{\color[HTML]{808080} \textit{RF-4.3}}}                & \multicolumn{5}{l|}{}                                                                                             \\ \hline
    \multicolumn{1}{|l|}{Precondición}       & \multicolumn{7}{p{12cm}|}{{\color[HTML]{808080} \textit{Las reservas estén disponibles y previamente hayan sido registradas en el historial.}}}                                                      \\ \hline
    \multicolumn{1}{|l|}{Poscondición}       & \multicolumn{7}{l|}{}                                                                                                                                                                          \\ \hline
    \multicolumn{1}{|l|}{Autor}              & \multicolumn{1}{l|}{}                         & \multicolumn{2}{l|}{Fecha}             & \multicolumn{1}{l|}{}       & \multicolumn{2}{l|}{Versión}           & 1.0                            \\ \hline
    \multicolumn{8}{|l|}{\cellcolor[HTML]{F2F2F2}Propósito}                                                                                                                                                                                   \\ \hline
    \multicolumn{8}{|p{15cm}|}{{\color[HTML]{808080} \textit{Proveer información detallada sobre destinos turísticos.}}}                                                                                                                            \\ \hline
    \multicolumn{8}{|l|}{\cellcolor[HTML]{F2F2F2}Resumen}                                                                                                                                                                                     \\ \hline
    \multicolumn{8}{|p{15cm}|}{{\color[HTML]{999999} \textit{El sistema proporciona información sobre destinos turísticos, incluyendo descripciones, imágenes y recomendaciones. Puede integrar IA para personalizar la experiencia del usuario.}}} \\ \hline
    \end{tabular}
    \end{table}

% Please add the following required packages to your document preamble:
% \usepackage[table,xcdraw]{xcolor}
% Beamer presentation requires \usepackage{colortbl} instead of \usepackage[table,xcdraw]{xcolor}
% Please add the following required packages to your document preamble:
% \usepackage[table,xcdraw]{xcolor}
% Beamer presentation requires \usepackage{colortbl} instead of \usepackage[table,xcdraw]{xcolor}
\begin{table}[H]
    \begin{tabular}{|llllllll|}
    \hline
    \multicolumn{1}{|l|}{Caso de Uso}                   & \multicolumn{5}{l|}{Gestionar cambios de reservas de forma autónoma}                                                                                                                     & \multicolumn{2}{l|}{\cellcolor[HTML]{E6E6E6}CU-34}                                 \\ \hline
    \multicolumn{1}{|l|}{Actores}                       & \multicolumn{7}{l|}{{\color[HTML]{808080} \textit{Cliente, Empleado}}}                                                                                                                                                                                                        \\ \hline
    \multicolumn{1}{|l|}{Tipo}                          & \multicolumn{7}{l|}{{\color[HTML]{808080} \textit{Primario, Esencial}}}                                                                                                                                                                                                       \\ \hline
    \multicolumn{1}{|l|}{Referencias}                   & \multicolumn{2}{l|}{{\color[HTML]{808080} \textit{RF-4.2}}}                                     & \multicolumn{5}{l|}{{\color[HTML]{808080} \textit{CU-33}}}                                                                                                                  \\ \hline
    \multicolumn{1}{|l|}{Precondición}                  & \multicolumn{7}{l|}{{\color[HTML]{808080} \textit{La presencia de IA en el sistema.}}}                                                                                                                                                                                        \\ \hline
    \multicolumn{1}{|l|}{Poscondición}                  & \multicolumn{7}{l|}{}                                                                                                                                                                                                                                                         \\ \hline
    \multicolumn{1}{|l|}{Autor}                         & \multicolumn{1}{l|}{}                                    & \multicolumn{2}{l|}{Fecha}                                   & \multicolumn{1}{l|}{}                    & \multicolumn{2}{l|}{Versión}                                 & 1.0                                       \\ \hline
    \multicolumn{8}{|l|}{\cellcolor[HTML]{F2F2F2}Propósito}                                                                                                                                                                                                                                                                             \\ \hline
    \multicolumn{8}{|p{15cm}|}{{\color[HTML]{808080} \textit{Permitir a los usuarios modificar sus reservas sin intervención humana.}}}                                                                                                                                                                                                       \\ \hline
    \multicolumn{8}{|l|}{\cellcolor[HTML]{F2F2F2}Resumen}                                                                                                                                                                                                                                                                               \\ \hline
    \multicolumn{8}{|p{15cm}|}{{\color[HTML]{999999} \textit{El sistema facilita la gestión de cambios en reservas de forma autónoma mediante IA, permitiendo modificaciones como cambios de fecha, cancelaciones o ajustes según disponibilidad y políticas establecidas. En casos excepcionales, se podrá escalar a intervención humana.}}} \\ \hline
    \end{tabular}
    \end{table}

% Please add the following required packages to your document preamble:
% \usepackage[table,xcdraw]{xcolor}
% Beamer presentation requires \usepackage{colortbl} instead of \usepackage[table,xcdraw]{xcolor}
\begin{table}[H]
    \begin{tabular}{|llllllll|}
    \hline
    \multicolumn{1}{|l|}{Caso de Uso}                    & \multicolumn{5}{l|}{Soporte y Asistencia}                                                                                                                                               & \multicolumn{2}{l|}{\cellcolor[HTML]{E6E6E6}CU-35}                                    \\ \hline
    \multicolumn{1}{|l|}{Actores}                        & \multicolumn{7}{l|}{{\color[HTML]{808080} \textit{Administrador, Cliente, Guía Turístico, Empleado}}}                                                                                                                                                                           \\ \hline
    \multicolumn{1}{|l|}{Tipo}                           & \multicolumn{7}{l|}{{\color[HTML]{808080} \textit{Primario, Esencial}}}                                                                                                                                                                                                         \\ \hline
    \multicolumn{1}{|l|}{Referencias}                    & \multicolumn{2}{l|}{{\color[HTML]{808080} \textit{RF-4}}}                                         & \multicolumn{5}{l|}{{\color[HTML]{808080} \textit{CU-29, CU-30}}}                                                                                                           \\ \hline
    \multicolumn{1}{|l|}{Precondición}                   & \multicolumn{7}{l|}{{\color[HTML]{808080} \textit{Disponibilidad del soporte con ayuda de IA.}}}                                                                                                                                                                                \\ \hline
    \multicolumn{1}{|l|}{Poscondición}                   & \multicolumn{7}{l|}{}                                                                                                                                                                                                                                                           \\ \hline
    \multicolumn{1}{|l|}{Autor}                          & \multicolumn{1}{l|}{}                                    & \multicolumn{2}{l|}{Fecha}                                    & \multicolumn{1}{l|}{}                  & \multicolumn{2}{l|}{Versión}                                   & 1.0                                        \\ \hline
    \multicolumn{8}{|l|}{\cellcolor[HTML]{F2F2F2}Propósito}                                                                                                                                                                                                                                                                                \\ \hline
    \multicolumn{8}{|p{15cm}|}{{\color[HTML]{808080} \textit{Brindar soporte y asistencia a los usuarios en diferentes aspectos del sistema.}}}                                                                                                                                                                                                  \\ \hline
    \multicolumn{8}{|l|}{\cellcolor[HTML]{F2F2F2}Resumen}                                                                                                                                                                                                                                                                                  \\ \hline
    \multicolumn{8}{|p{15cm}|}{{\color[HTML]{999999} \textit{El sistema proporciona soporte y asistencia a los usuarios mediante diferentes canales, incluyendo chat en vivo, preguntas frecuentes y contacto directo con un agente cuando sea necesario. Puede integrar IA para resolver consultas comunes y optimizar tiempos de respuesta.}}} \\ \hline
    \end{tabular}
    \end{table}

    % Please add the following required packages to your document preamble:
% \usepackage[table,xcdraw]{xcolor}
% Beamer presentation requires \usepackage{colortbl} instead of \usepackage[table,xcdraw]{xcolor}
\begin{table}[H]
    \begin{tabular}{|llllllll|}
    \hline
    \multicolumn{1}{|l|}{Caso de Uso}       & \multicolumn{5}{l|}{Identificación de Turistas}                                                                          & \multicolumn{2}{l|}{\cellcolor[HTML]{E6E6E6}CU-36}          \\ \hline
    \multicolumn{1}{|l|}{Actores}           & \multicolumn{7}{l|}{{\color[HTML]{808080} \textit{Cliente}}}                                                                                                                           \\ \hline
    \multicolumn{1}{|l|}{Tipo}              & \multicolumn{7}{l|}{{\color[HTML]{808080} \textit{Primario, Esencial}}}                                                                                                                \\ \hline
    \multicolumn{1}{|l|}{Referencias}       & \multicolumn{2}{l|}{{\color[HTML]{808080} \textit{RF-4.1}}}              & \multicolumn{5}{l|}{}                                                                                       \\ \hline
    \multicolumn{1}{|l|}{Precondición}      & \multicolumn{7}{l|}{{\color[HTML]{808080} \textit{Tenga la opción disponible en web.}}}                                                                                                \\ \hline
    \multicolumn{1}{|l|}{Poscondición}      & \multicolumn{7}{l|}{}                                                                                                                                                                  \\ \hline
    \multicolumn{1}{|l|}{Autor}             & \multicolumn{1}{l|}{}                        & \multicolumn{2}{l|}{Fecha}           & \multicolumn{1}{l|}{}     & \multicolumn{2}{l|}{Versión}         & 1.0                           \\ \hline
    \multicolumn{8}{|l|}{\cellcolor[HTML]{F2F2F2}Propósito}                                                                                                                                                                          \\ \hline
    \multicolumn{8}{|p{15cm}|}{{\color[HTML]{808080} \textit{Reservar y comprar paquetes de manera online a través de la web.}}}                                                                                                           \\ \hline
    \multicolumn{8}{|l|}{\cellcolor[HTML]{F2F2F2}Resumen}                                                                                                                                                                            \\ \hline
    \multicolumn{8}{|p{15cm}|}{{\color[HTML]{808080} \textit{Se permite reservar y comprar paquetes de manera online por parte del cliente gracias a la implementación de la web por parte de los trabajadores de la empresa turística.}}} \\ \hline
    \end{tabular}
    \end{table}

